\section{Cote 66, subject 21: the exact affine henselian comparison}

\subsection{Original statement and coefficient scope}
The official cote 66 PDF, physical page 4, archive footer 3, subject 21,
asks when
\[
 H^*(X,F)\longrightarrow H^*(X_0,F_0)
\]
is an isomorphism for a closed subscheme $X_0\subset X$, and asks
\emph{``Suffit-il que ce soit vrai en dim 0?''}
The margin identifies \emph{cohomologie \'etale}. The specified cases
are $X=\operatorname{Spec}A$, $X_0=\operatorname{Spec}(A/I)$ with
$A$ separated and complete for the $I$-adic topology, and \'etale
localization along a closed subscheme of an affine scheme.
The sheet does not explicitly write the quantifier on $F$ or restrict
it to torsion. We state those choices: the all-torsion-sheaf affine
question has the following exact answer, while a test on one fixed
sheaf does not suffice.

\subsection{BR-66-01. Degree zero for all torsion coefficients is an exact affine test}
Let $A$ be a commutative ring and $I\subset A$ any ideal, and set
$X=\operatorname{Spec}A$, $Z=\operatorname{Spec}(A/I)$. The following
conditions are equivalent:
\begin{enumerate}
\item $(A,I)$ is a henselian pair.
\item For every torsion abelian sheaf $F$ on $X_{\mathrm{\acute et}}$,
      the restriction $H^0(X,F)\to H^0(Z,F|_Z)$ is an isomorphism.
\item For every such $F$ and every $q\geq0$,
      $H^q(X,F)\to H^q(Z,F|_Z)$ is an isomorphism.
\item For every finite morphism $f:Y\to X$, restriction gives a
      bijection from the clopen subsets of $Y$ to those of
      $Y\times_X Z$.
\end{enumerate}
In particular, the degree-zero condition can be tested just on the
sheaves $f_*\underline{\mathbf F}_2$ for all finite $f$.

\paragraph{Exact known inputs.}
Gabber's affine analogue of proper base change gives (1)$\Rightarrow$(3)
for arbitrary torsion sheaves, with no noetherian assumption and no
invertibility assumption on their torsion orders. The exact statement
used is The Stacks Project, Theorem 59.82.7, Tag \texttt{09ZI}; the
original article is Gabber (1994), pp.~325--335. Its complete published
proof is not reproduced here. The equivalence (1)$\Leftrightarrow$(4)
is the idempotent criterion in The Stacks Project, Lemma 15.11.6(1),(3),
Tag \texttt{09XI}. We prove the short implication from the stated
sheaf test to that criterion, to settle the quantifier precisely.

\paragraph{Complete proof of the additional implication and equivalence.}
The implication (3)$\Rightarrow$(2) is immediate. Assume (2), and
fix a finite morphism $f:Y\to X$. The sheaf
$F=f_*\underline{\mathbf F}_2$ is annihilated by $2$, so (2) applies.
Put $Y_Z=Y\times_X Z$ and $f_Z:Y_Z\to Z$. Proper base change in
degree zero for the finite morphism gives
\[
 (f_*\underline{\mathbf F}_2)|_Z
       \simeq (f_Z)_*\underline{\mathbf F}_2.
\]
An exact inspected locator for this known input is The Stacks
Project, Theorem 59.91.11, Tag \texttt{095T}; only its degree-zero
finite-morphism case is used. By the definition of direct image,
condition (2) therefore gives
\[
 \Gamma(Y,\underline{\mathbf F}_2)
       \xrightarrow{\sim}\Gamma(Y_Z,\underline{\mathbf F}_2).
 \tag{HEN.1}
\]
A section of the constant two-element \'etale sheaf is the same as
a decomposition of the scheme into two open and closed subsets:
it is a section of the disjoint union of two copies of the scheme,
and the inverse images of its two components are the two subsets.
Thus (HEN.1) is precisely condition (4). Since $Y$ is affine,
these decompositions are also precisely the idempotents of its ring,
and restriction is reduction modulo $I$. The cited idempotent
criterion implies (1). The same proof works if (2) is assumed only
for this displayed family of finite direct images. Combining it with
Gabber's theorem proves all stated equivalences. \(\square\)

\subsection{BR-66-02. Both affine cases named on the sheet are affirmative}
If $A$ is $I$-adically separated and complete, the desired comparison
holds for every torsion abelian \'etale sheaf and every degree.
The same holds for the henselization $(A^h,IA^h)$ of any affine pair,
with closed subscheme $\operatorname{Spec}(A/I)$.
These are direct applications of known results: The Stacks Project,
Lemma 15.11.4, Tag \texttt{0ALJ}, says an $I$-adically complete pair is
henselian; Lemma 15.12.1, Tag \texttt{0A02}, constructs the henselization
as the filtered colimit of \'etale neighborhoods inducing an isomorphism
on the closed subscheme, proves it henselian, and identifies
$A^h/IA^h=A/I$. Apply BR-66-01. This identifies the affine \'etale
localization intended by the second case without asserting that an
arbitrary unrelated pair is henselian.

\subsection{BR-66-03. Testing only one coefficient sheaf in degree zero fails}
Even for smooth affine schemes over $\mathbf C$, comparison in degree
zero for one fixed finite constant sheaf does not imply comparison
in degree one for that same sheaf.

\paragraph{Complete counterexample.}
Take $X=\operatorname{Spec}\mathbf C[t,t^{-1}]$,
$Z=\{t=1\}=\operatorname{Spec}\mathbf C$, and
$F=\underline{\mathbf Z/2}$. Both schemes are connected, so
restriction on $H^0$ is the identity on $\mathbf Z/2$.
The morphism
\[
 \operatorname{Spec}\mathbf C[u,u^{-1}]\longrightarrow X,
 \qquad t=u^2,
\]
is a torsor for $\{1,-1\}\simeq\mathbf Z/2$. It is finite \'etale:
its ring is free of rank two over $\mathbf C[t,t^{-1}]$, and the
derivative $2u$ is a unit. The sign action is free and acts transitively
on every geometric fibre. It is not a trivial torsor, because a section
would give a unit $a(t)\in\mathbf C[t,t^{-1}]$ with $a(t)^2=t$.
Every Laurent-polynomial unit is $ct^m$: the least and greatest
exponents of a product are sums of the corresponding exponents,
so a product equal to $1$ forces the difference between the greatest
and least exponent of each factor to be zero. Then $a(t)^2=t$
would require $2m=1$, which is impossible.

The classification of abelian-sheaf torsors by $H^1$ (The Stacks
Project, Lemma 21.4.3, Tag \texttt{03AJ}, with the section criterion
in Lemma 21.4.2, Tag \texttt{03AI}) therefore gives
a nonzero class in $H^1(X,F)$. Its restriction to $Z$ is trivial:
the fibre consists of the two points $u=1,-1$ and has a section.
Consequently $H^1(X,F)\to H^1(Z,F|_Z)$ is not injective.
This proves the assertion. \(\square\)

The result answers the original affine torsion formulation, including
its degree-zero test with the universal coefficient quantifier.
It does not assert a general comparison theorem for arbitrary
quasicompact schemes or for nontorsion coefficients. The unseen
SA22 run's specific nonseparated counterexample and its residual
\texttt{U:henselian-comparison} are not identified with this elementary
fixed-sheaf counterexample.

\subsection*{Exact references and source inspection}
Grothendieck, A. (n.d.). \emph{List of subjects}, cote 66,
physical PDF p.~4, archive footer 3, subject 21.
\url{https://grothendieck.umontpellier.fr/66.pdf}.
The complete sheet and an enlarged subject-21 crop were visually read
in this run; ``66/004'' in the intake denotes the physical PDF page.

Gabber, O. (1994). Affine analog of the proper base change theorem.
\emph{Israel Journal of Mathematics, 87}, 325--335.
\url{https://doi.org/10.1007/BF02773001}.
The publisher abstract and bibliographic record were inspected;
the accessible complete theorem and proof used are the following
Stacks Project version, so no unretrieved theorem number from the
original article is asserted as newly checked.

The Stacks Project Authors. (2026). \emph{The Stacks Project}.
Theorem 59.82.7 (Gabber), \url{https://stacks.math.columbia.edu/tag/09ZI};
Lemma 15.11.6, \url{https://stacks.math.columbia.edu/tag/09XI};
Theorem 59.91.11, \url{https://stacks.math.columbia.edu/tag/095T};
Lemma 15.11.4, \url{https://stacks.math.columbia.edu/tag/0ALJ};
Lemma 15.12.1, \url{https://stacks.math.columbia.edu/tag/0A02}.
Lemmas 21.4.2--21.4.3 in \url{https://stacks.math.columbia.edu/tag/03AG}.
All displayed statements were directly inspected on 21 November 2026.
